% 依据原卷第1页 q19 图：两组医学指标频率/组距直方图。
% 患病者分组 [95,100),[100,105),...,[125,130) 的频率组距分别为
% 0.002, 0.012, 0.034, 0.036, 0.040, 0.040, 0.036。
% 未患病者分组 [70,75),[75,80),...,[100,105) 的频率组距分别为
% 0.038, 0.040, 0.040, 0.036, 0.034, 0.010, 0.002。
\begin{tikzpicture}[font=\normalsize]
  \begin{axis}[exam statistical histogram,
    width=6.55cm,
    height=13.689cm,
    xmin=90, xmax=132,
    ymin=0, ymax=0.045,
    axis lines=left,
    axis line style={-latex},
    tick style={black},
    xtick={95,100,105,110,115,120,125,130},
    ytick={0.002,0.012,0.034,0.036,0.040},
    yticklabels={0.002,0.012,0.034,0.036,0.040},
    scaled y ticks=false,
    xticklabel style={font=\normalsize},
    yticklabel style={font=\normalsize},
    tick align=outside,
    clip=false,
  ]
    % Each ybar interval is one width-5 class; bars are unfilled outlines.
    \addplot+[ybar interval, mark=none, draw=black, fill=white, line width=0.35pt]
      coordinates {(95,0.002) (100,0.012) (105,0.034) (110,0.036)
                   (115,0.040) (120,0.040) (125,0.036) (130,0.036)};
    % Dashed guide segments terminate at the corresponding bar boundary.
    \draw[dashed] (axis cs:90,0.002) -- (axis cs:95,0.002);
    \draw[dashed] (axis cs:90,0.012) -- (axis cs:105,0.012);
    \draw[dashed] (axis cs:90,0.034) -- (axis cs:105,0.034);
    \draw[dashed] (axis cs:90,0.036) -- (axis cs:110,0.036);
    \draw[dashed] (axis cs:90,0.040) -- (axis cs:115,0.040);
    % Broken origin on the compressed x-axis.
    \draw (axis cs:91.0,0) -- (axis cs:91.5,0.0008) -- (axis cs:92.0,0)
          -- (axis cs:92.5,0.0008) -- (axis cs:93.0,0);
    \node[anchor=north east, inner sep=1pt, font=\normalsize,yshift=-4pt] at (axis cs:90,0.000000) {0};
    \node[anchor=north, inner sep=1pt] at (axis description cs:0.52,-0.27) {患病者};
    \examstatxlabel{指标};
    \examstatylabel{\(\dfrac{\text{频率}}{\text{组距}}\)};
  \end{axis}

  \begin{scope}[xshift=7.15cm]
    \begin{axis}[exam statistical histogram,
      width=6.55cm,
      height=13.689cm,
      xmin=65, xmax=107,
      ymin=0, ymax=0.045,
      axis lines=left,
      axis line style={-latex},
      tick style={black},
      xtick={70,75,80,85,90,95,100,105},
      ytick={0.002,0.010,0.034,0.036,0.038,0.040},
      yticklabels={0.002,0.010,0.034,0.036,0.038,0.040},
      scaled y ticks=false,
      xticklabel style={font=\normalsize},
      yticklabel style={font=\normalsize},
      tick align=outside,
      clip=false,
    ]
      \addplot+[ybar interval, mark=none, draw=black, fill=white, line width=0.35pt]
        coordinates {(70,0.038) (75,0.040) (80,0.040) (85,0.036)
                     (90,0.034) (95,0.010) (100,0.002) (105,0.002)};
      \draw[dashed] (axis cs:65,0.002) -- (axis cs:70,0.002);
      \draw[dashed] (axis cs:65,0.010) -- (axis cs:100,0.010);
      \draw[dashed] (axis cs:65,0.034) -- (axis cs:95,0.034);
      \draw[dashed] (axis cs:65,0.036) -- (axis cs:90,0.036);
      \draw[dashed] (axis cs:65,0.038) -- (axis cs:75,0.038);
      \draw[dashed] (axis cs:65,0.040) -- (axis cs:80,0.040);
      \draw (axis cs:66.0,0) -- (axis cs:66.5,0.0008) -- (axis cs:67.0,0)
            -- (axis cs:67.5,0.0008) -- (axis cs:68.0,0);
      \node[anchor=north east, inner sep=1pt, font=\normalsize,yshift=-4pt] at (axis cs:65,0.000000) {0};
      \node[anchor=north, inner sep=1pt] at (axis description cs:0.52,-0.27) {未患病者};
      \examstatxlabel{指标};
      \examstatylabel{\(\dfrac{\text{频率}}{\text{组距}}\)};
    \end{axis}
  \end{scope}
\end{tikzpicture}
